A popularização das balanças inteligentes de bioimpedância ampliou o acesso a
informações sobre composição corporal, como percentual de gordura, massa muscular,
metabolismo basal e gordura visceral. Mas a maioria dos aplicativos só mostra números e
classificações genéricas, e exige do usuário algum conhecimento técnico para interpretar
os resultados.

Além disso, as medições por bioimpedância são influenciadas por diversos fatores
fisiológicos e comportamentais, como hidratação, alimentação recente, prática de
exercícios físicos e horário da avaliação. Tais fatores podem introduzir variações
significativas nos resultados, dificultando a interpretação da evolução corporal ao longo
do tempo.

Falta, portanto, uma ferramenta que não apenas registre os dados das balanças, mas também
os contextualize e interprete de forma personalizada. Juntar bioimpedância, inteligência
artificial e tecnologias móveis permite transformar dados brutos em informações úteis para
a saúde.

O \sintonia\ foi concebido com esse propósito. Além de se integrar à balança CF577BLE,
reúne um assistente virtual baseado em IA, o registro de contexto por anamnese e um modelo
computacional que contextualiza e ajusta a interpretação das medições, para dar ao usuário
uma visão mais confiável e personalizada da composição corporal. Oferecer a interação por
aplicativo web progressivo (PWA) e por WhatsApp amplia o acesso à solução e facilita o
acompanhamento contínuo e o autocuidado baseado em dados.
